
První případ užití je specifikován v tabulce č. \ref{tab:usecase:UC01}.

\begin{UseCaseTable}
  {Registrace do systému}
  {UC01}
  {F01}
  {Vytvoření účtu provozovatele monitorovaného prostoru za účelem evidence a publikace měření}
  {Provozovatel monitorovaného prostoru}
  {Emailová služba}
  {Uživatel není v systému evidován}
  {Uživatel je v systému evidován v neaktivním stavu, byl odeslán potvrzovací email s odkazem pro aktivaci účtu}

\UseCaseScenario{Základní scénář}
\UseCaseStep{uživatel}{Otevře domovskou stránku a klikne na tlačítko pro registraci.}
\UseCaseStep{systém}{Vrátí registrační formulář.}
\UseCaseStep{uživatel}{Vyplní email a heslo a formulář odešle.}
\UseCaseStep{systém}{Ověří, že zadaný email dosud není registrován.}
\UseCaseStep{systém}{Vytvoří uživatelský účet v neaktivním stavu.}
\UseCaseStep{systém}{Odešle na zadaný email zprávu s odkazem pro aktivaci účtu.}

\UseCaseScenario{Alternativní scénář -- email již registrován}
\UseCaseStepCustom{A1}{uživatel, systém}{shodné s kroky 1--3 základního scénáře.}
\UseCaseStepCustom{A2}{systém}{Zjistí, že zadaný email je již registrován, a vyzve uživatele k zadání jiného emailu.}
\UseCaseStepCustom{A3}{uživatel}{Zadá jiný email.}
\UseCaseStepCustom{A4}{systém}{Pokračuje krokem 4 základního scénáře.}

\end{UseCaseTable}
